% sections/intro.tex
% Seção 1: Introdução - O que é o IntelliStock

\section{Introdução}

\subsection{O Desafio do Estoque em Food Service}

Restaurantes operam em um equilíbrio delicado: ter estoque suficiente para
atender a demanda sem gerar desperdício de produtos perecíveis.

Os impactos de um controle ineficiente são diretos:

\begin{itemize}[leftmargin=2em]
  \item \textbf{Falta de ingredientes} durante horários de pico, resultando em
        pratos indisponíveis e perda de receita
  \item \textbf{Desperdício} de produtos que vencem antes do uso
  \item \textbf{Decisões reativas} tomadas sob pressão, sem visão do cenário
        completo
  \item \textbf{Dependência} da experiência individual de funcionários
\end{itemize}

\subsection{Por que Sistemas Tradicionais Não Resolvem}

Sistemas ERP e PDV convencionais oferecem controle de estoque básico:

\begin{itemize}[leftmargin=2em]
  \item Registro de entrada e saída
  \item Relatórios de movimentação
  \item Alertas de estoque mínimo fixo
\end{itemize}

Porém, \textbf{não entregam inteligência} sobre:

\begin{itemize}[leftmargin=2em]
  \item Demanda futura baseada em histórico de vendas
  \item Tempo de entrega (lead time) de cada fornecedor
  \item Validade e risco de perda de produtos em estoque
  \item Necessidade de produção de itens semi-acabados
  \item Explosão automática de receitas (BOM)
\end{itemize}

\subsection{IntelliStock: Inteligência para Decisões de Estoque}

O IntelliStock é um \textbf{motor de inteligência de estoque} projetado
especificamente para operações de food service.

Ele analisa o estado atual do estoque, a demanda histórica e prevista,
e gera \textbf{alertas acionáveis} para três perfis:

\begin{center}
\begin{tabular}{>{\bfseries}l p{8cm}}
\toprule
Perfil & Alertas \\
\midrule
Compras   & O que comprar, quando, quanto \\
Cozinha   & O que produzir, prioridades de produção \\
Gestão    & Riscos de desperdício, produtos parados, visão geral \\
\bottomrule
\end{tabular}
\end{center}

\begin{infobox}[Conceito Central]
O IntelliStock \textbf{complementa} sistemas existentes. Ele recebe dados do
ERP/PDV atual e retorna recomendações inteligentes de ação.
\end{infobox}

\subsection{Principais Capacidades}

O IntelliStock entrega:

\begin{itemize}[leftmargin=2em]
  \item Pontos de reposição dinâmicos baseados em demanda real
  \item Consideração automática de lead time de fornecedores
  \item Explosão de receitas para antecipar necessidade de insumos
  \item Gestão FEFO (First Expired, First Out) para controle de validade
  \item Diferenciação entre urgências e planejamento
  \item Classificação ABC de criticidade de itens
  \item Alertas direcionados por responsabilidade
\end{itemize}

\subsection{Escopo do Produto}

\begin{infobox}[Modelo de Atuação]
O IntelliStock atua como \textbf{camada de inteligência}, não como substituto
de sistemas operacionais.
\end{infobox}

\textbf{O IntelliStock é:}
\begin{itemize}[leftmargin=2em]
  \item Motor de análise e recomendação
  \item Complemento ao ERP/PDV existente
  \item Ferramenta de apoio à decisão
\end{itemize}

\textbf{O IntelliStock não é:}
\begin{itemize}[leftmargin=2em]
  \item Sistema de PDV ou emissão de notas
  \item ERP com controle de entrada/saída
  \item Sistema de compras automáticas
\end{itemize}

\subsection{Público-Alvo}

O IntelliStock foi projetado para operações de food service que:

\begin{itemize}[leftmargin=2em]
  \item Trabalham com produtos perecíveis
  \item Possuem produção interna (cozinha que transforma insumos)
  \item Precisam antecipar demanda com base em histórico
  \item Buscam reduzir desperdício sem comprometer atendimento
  \item Já utilizam ERP/PDV e desejam inteligência adicional
\end{itemize}
